\section{Introduction}
A point on the ordinary real line carries no internal left--right multiplicity: the coordinate $x$ is a single point. The split or double-arrow construction replaces $x$ by two immediately adjacent points $x^-<x^+$ with no point strictly between them, while a projection $\pi$ forgets the distinction and sends both sides to $x$. As a topological construction this is classical. Its relevance here is not that quantum mechanics secretly lives on a double-arrow space, but that the split topology supplies a precise way to make a two-sided jump continuous before a coarse projection identifies the two sides.

The motivating operator-theoretic fact is equally elementary. For noncommuting observables $A$ and $B$, the products $AB$ and $BA$ are distinct. Their Jordan and Lie components, $(AB+BA)/2$ and $[A,B]/(2i)$, are standard. Merely renaming these parts would add no mathematics. The question pursued here is instead whether the order pair can be lifted to a genuine continuous object on a split base in such a way that noncommutativity acquires a universal descent meaning, and whether that local construction extends nontrivially to the space of all orderings of a finite operator family.

The answer is affirmative. The local lift turns $[A,B]$ into jump data and $\|[A,B]\|/2$ into an exact distance to the subspace of sections that forget the split. For unitary families, the local weights induce a bottleneck geometry on the Cayley graph of $S_n$. This geometry admits a closed inversion-set formula, a threshold filtration by weighted incompatibility graphs, and two oriented enrichments. One of those enrichments reveals a sharp distinction between full operator context and pair-only reduction: the exact contextual connection is flat, while the reduced pair connection can have nontrivial holonomy.

No new physical postulate is introduced. The paper remains within standard $C^*$-algebraic operator mechanics. The contribution is a representation and a family of derived invariants. In particular, the words \emph{curvature} and \emph{holonomy} are used only for a discrete operator-order connection defined below; no spacetime, gauge-field, or dynamical interpretation is inferred.

\paragraph{Scope and proof status.} All statements labeled Theorem, Proposition, or Corollary are proved here or explicitly marked as known infrastructure. Conjectural spectral and Clifford extensions are isolated in \cref{sec:related} and are not used as premises for proved results.

\paragraph{Main proved contributions.}
\begin{itemize}
\item An operator order-lift $\Ot(A,B)$ on a split interval, with commutativity equivalent to descent to the ordinary interval.
\item The exact obstruction formula $\dist(\Ot(A,B),\CI(\A))=\|[A,B]\|/2$.
\item A minimax pseudoultrametric on the set of operator orderings and, for unitary families, a closed inversion-set formula.
\item A threshold description of order clusters via acyclic orientations of weighted incompatibility graphs, plus exact reconstruction of the pairwise commutator-norm matrix from the marked order geometry.
\item Parallel additive and multiplicative oriented lifts: a Lie-commutator cocycle and a group-commutator transport.
\item Exact flatness of full contextual transport, together with generally nontrivial pair-reduced square and braid holonomies.
\item A near-identity hierarchy separating pairwise commutators, Jacobi-type context corrections, and fourth-order pair-reduced curvature.
\item Closed spin-$\tfrac12$ formulas and an exact nontrivial braid-holonomy example.
\item Exact fixed-target and precedence-target unitary-circuit reordering certificates, including simultaneous optimality of reduced paths for all pair-weighted $\ell_p$ swap costs.
\end{itemize}

The original motivating examples have different statuses. The operator-order lift and the spin-$\tfrac12$ order geometry are proved here. The Robertson--Schr\"odinger relation and the Pauli magnetic term are retained as standard state-level and physical illustrations, respectively, rather than novelty claims. The resolvent-boundary construction remains an explicit conjectural direction.

\section{Split intervals and descent spaces}
Let $I=[a,b]$. Define the compact split interval
\[
I_s=\{a^+\}\cup\{t^-,t^+:a<t<b\}\cup\{b^-\},
\]
ordered lexicographically so that $t^-<t^+$ and no point lies strictly between them. The projection $\pi:I_s\to I$ is defined by $\pi(t^\pm)=t$ and the evident endpoint convention.

For a Banach space $E$, set
\[
\SI(E):=C(I_s,E),\qquad \CI(E):=\pi^*C(I,E),
\]
with the supremum norm. Thus $F\in\CI(E)$ exactly when $F(t^-)=F(t^+)$ at every split point. Write
\[
\Delta_tF:=F(t^+)-F(t^-).
\]
The pullback $\pi^*:C(I,E)\to C(I_s,E)$ is isometric because $\pi$ is surjective; in particular, $\CI(E)$ is closed.

\begin{theorem}[Distance to the descent subspace; known split/regulated-function infrastructure]\label{thm:distance}
For every $F\in C(I_s,E)$,
\begin{equation}\label{eq:distance}
\dist\bigl(F,\CI(E)\bigr)=\frac12\sup_{t\in(a,b)}\|\Delta_tF\|.
\end{equation}
\end{theorem}
\begin{proof}
Set $M(F)=\sup_t\|\Delta_tF\|$. If $G=g\circ\pi\in\CI(E)$, then $G(t^-)=G(t^+)=g(t)$ and hence
\[
\|\Delta_tF\|\le \|F(t^+)-g(t)\|+\|g(t)-F(t^-)\|\le 2\|F-G\|_\infty.
\]
Thus $\dist(F,\CI(E))\ge M(F)/2$.

For the reverse inequality, fix $\varepsilon>0$. Under the standard split-line/regulated-function correspondence, $F$ corresponds to a Banach-valued regulated function on $I$; regulated functions are uniform limits of finite step functions, and lifting the one-sided step values gives finite clopen step sections on $I_s$ \cite{ref1,ref14}. Choose such an $S$ with $\|F-S\|_\infty<\varepsilon$. Let $v_1,\dots,v_m$ be its consecutive plateau values and put $M(S)=\max_k\|v_{k+1}-v_k\|$. Then $M(S)\le M(F)+2\varepsilon$. At each interior jump set $m_k=(v_k+v_{k+1})/2$. On every interior plateau with value $v_k$, interpolate affinely between its two neighboring midpoint values. On the first plateau interpolate from $g(a)=v_1$ to $m_1$, and on the last interpolate from $m_{m-1}$ to $g(b)=v_m$ (with the constant choice if $m=1$). In each case the interpolation endpoints lie in the closed ball $\overline B(v_k,M(S)/2)$ around the relevant plateau value; norm balls in a Banach space are convex, so the entire affine segment remains in that ball. This produces $g\in C(I,E)$ satisfying
\[
\|S-g\circ\pi\|_\infty\le \frac12M(S).
\]
Therefore
\[
\dist(F,\CI(E))\le \varepsilon+\frac12M(S)\le \frac12M(F)+2\varepsilon.
\]
Letting $\varepsilon\downarrow0$ proves \eqref{eq:distance}.
\end{proof}

The functional-analytic infrastructure behind \cref{thm:distance} is classical and is not claimed as a new theorem. Continuous functions on split ordered spaces are closely tied to regulated and c\`adl\`ag data \cite{ref1,ref2,ref14}. Its role here is to provide a normed topological carrier for the operator-order lift.

\section{Operator order-lift and noncommutativity as descent obstruction}
Let $\A$ be a unital $C^*$-algebra and let $A,B\in\A$ be self-adjoint. Fix an interior split point $t$. Because $[a^+,t^-]$ and $[t^+,b^-]$ are clopen, the section
\begin{equation}\label{eq:orderlift}
\Ot(A,B)(s)=
\begin{cases}
BA,& s\le t^-,\\
AB,& s\ge t^+
\end{cases}
\end{equation}
is continuous on $I_s$. Its unique nonzero split jump is
\[
\Delta_t\Ot(A,B)=AB-BA=[A,B].
\]
The standard Jordan and Lie parts are recovered as
\begin{equation}\label{eq:jordanlie}
\mathcal C_t(\Ot)=\frac{AB+BA}{2},\qquad
\mathcal K_t(\Ot)=\frac{[A,B]}{2i}.
\end{equation}

\begin{theorem}[Commutativity--descent equivalence]\label{thm:commdescent}
For self-adjoint $A,B\in\A$, the following are equivalent:
\begin{enumerate}[label=(\roman*)]
\item $[A,B]=0$;
\item $\Ot(A,B)\in\CI(\A)$;
\item there exists $G\in C(I,\A)$ with $\Ot(A,B)=G\circ\pi$.
\end{enumerate}
\end{theorem}
\begin{proof}
If $[A,B]=0$, the lifted section is constant. Conversely, any factorization through $\pi$ forces equal values at $t^-$ and $t^+$, hence $BA=AB$.
\end{proof}

\begin{corollary}[Exact descent obstruction]\label{cor:descent}
\begin{equation}\label{eq:descent}
\desc(A,B):=\dist\bigl(\Ot(A,B),\CI(\A)\bigr)=\frac12\|[A,B]\|.
\end{equation}
\end{corollary}
\begin{proof}
Apply \cref{thm:distance} to the single jump $\Delta_t\Ot=[A,B]$.
\end{proof}

The phrase ``noncommutativity as an obstruction to descent'' will always mean the precise equivalence and distance formula above. It is not a new definition of the commutator and does not assert that all notions of quantum incompatibility reduce to this norm.

\section{Functoriality and representations}
A $*$-homomorphism $\varphi:\A\to\Balg$ induces
\[
\varphi_*:C(I_s,\A)\to C(I_s,\Balg),\qquad F\mapsto \varphi\circ F.
\]
It preserves the descent subspace and commutes with the order-lift:
\[
\varphi_*\Ot(A,B)=\Ot(\varphi(A),\varphi(B)).
\]

\begin{theorem}[Contractivity and faithful invariance]\label{thm:functor}
\[
\desc^{\Balg}(\varphi(A),\varphi(B))\le \desc^{\A}(A,B).
\]
If $\varphi$ is faithful, equality holds.
\end{theorem}
\begin{proof}
Every $*$-homomorphism of $C^*$-algebras is contractive, while every faithful $*$-homomorphism is isometric. Apply \eqref{eq:descent} to $\varphi([A,B])$.
\end{proof}

Thus the obstruction is intrinsic to the abstract $C^*$-algebraic pair and does not depend on a chosen faithful Hilbert-space realization.

\section{State evaluation and standard quantum geometry}
For a state $\rho$ with GNS representation $\pi_\rho$, define
\[
d_\rho(A,B)=\frac12\|\pi_\rho([A,B])\|,
\qquad
\omega_\rho(A,B)=\frac{1}{2i}\rho([A,B]).
\]
Then
\begin{equation}\label{eq:statehier}
|\omega_\rho(A,B)|\le d_\rho(A,B)\le \desc(A,B).
\end{equation}
A state therefore reads a scalar component of an operator-level obstruction that exists before the state is selected.

For centered observables, the real covariance coordinate and the commutator coordinate recover the standard Robertson--Schr\"odinger inequality. That inequality is not claimed as new. The distinction from geometric quantum mechanics is one of level: Kibble, Heslot, Ashtekar--Schilling, and Brody--Hughston develop geometry on the space of states or rays \cite{ref3,ref4,ref5,ref6}, whereas the split construction is defined at the operator-order level before state evaluation. The two descriptions are compatible rather than competing.

\subsection{Standard Pauli magnetic-term illustration}
Let $\Pi=p-qA$ be the kinetic momentum of a charged particle. Using $\sigma_i\sigma_j=\delta_{ij}I+i\varepsilon_{ijk}\sigma_k$ and
\[
[\Pi_i,\Pi_j]=iq\hbar\,\varepsilon_{ijk}B_k,
\]
one obtains
\begin{equation}\label{eq:pauli}
(\boldsymbol\sigma\cdot\boldsymbol\Pi)^2
=\boldsymbol\Pi^2+\frac{i}{2}\varepsilon_{ijk}\sigma_k[\Pi_i,\Pi_j]
=\boldsymbol\Pi^2-q\hbar\,\boldsymbol\sigma\cdot\mathbf B.
\end{equation}
Hence
\[
H_P=\frac{\boldsymbol\Pi^2}{2m}+q\phi-\frac{q\hbar}{2m}\boldsymbol\sigma\cdot\mathbf B.
\]
This is a standard physical illustration \cite{ref16}, not a new magnetic interaction.

\section{Finite operator-order geometry}
Let $A_1,\dots,A_n\in\A$ and define
\[
W_\sigma=A_{\sigma(1)}\cdots A_{\sigma(n)},\qquad \sigma\in S_n.
\]
Throughout the remainder of the paper, $[X,Y]=XY-YX$. If
\[
W_\sigma=P A_iA_jQ
\]
and $\sigma'$ swaps only that adjacent pair, then
\begin{equation}\label{eq:contexts}
W_{\sigma'}-W_\sigma=-P[A_i,A_j]Q,
\qquad
w_A(\sigma,\sigma'):=\frac12\|W_{\sigma'}-W_\sigma\|
=\frac12\|P[A_i,A_j]Q\|.
\end{equation}
A path of adjacent transpositions defines a finite split step-section whose descent distance is its largest edge weight. This motivates
\begin{equation}\label{eq:DA}
D_A(\sigma,\tau)=\min_{\gamma:\sigma\leadsto\tau}\max_{e\in\gamma}w_A(e).
\end{equation}

\begin{theorem}[Pseudoultrametric property; standard bottleneck-graph fact]\label{thm:pseudo}
For all $\sigma,\tau,\upsilon\in S_n$,
\[
D_A(\sigma,\sigma)=0,\qquad
D_A(\sigma,\tau)=D_A(\tau,\sigma),
\]
and
\[
D_A(\sigma,\upsilon)\le \max\{D_A(\sigma,\tau),D_A(\tau,\upsilon)\}.
\]
\end{theorem}
\begin{proof}
Reverse paths have the same weights. Concatenating a minimax path from $\sigma$ to $\tau$ with one from $\tau$ to $\upsilon$ produces a path whose bottleneck is the maximum of the two bottlenecks; minimization can only decrease it.
\end{proof}

\section{Unitary families: closed inversion formula}
Assume $U_1,\dots,U_n$ are unitary. Prefix and suffix factors in a neighboring swap are then unitary, so contextual multiplication does not change the $C^*$-norm. The edge cost becomes pair-intrinsic:
\begin{equation}\label{eq:cij}
c_{ij}=\frac12\|[U_i,U_j]\|.
\end{equation}
For $\sigma,\tau\in S_n$, let $\Inv(\sigma,\tau)$ be the unordered pairs whose relative order differs between $\sigma$ and $\tau$.

\begin{theorem}[Inversion-set formula]\label{thm:inversion}
\begin{equation}\label{eq:inversion}
D_U(\sigma,\tau)=\max\{c_{ij}:\{i,j\}\in\Inv(\sigma,\tau)\},
\end{equation}
with the maximum over the empty set equal to zero.
\end{theorem}
\begin{proof}
Every inverted pair must be swapped directly at least once along any adjacent-transposition path, which gives the lower bound. A reduced adjacent-transposition sorting path changes only inverted pairs, giving the matching upper bound.
\end{proof}

After the $O(n^2)$ pairwise commutator data have been computed, a distance query reduces to an inversion-set maximum. This differs from weighted Kendall metrics, where adjacent-swap costs are accumulated additively along shortest paths \cite{ref7,ref17}.

For normalized spin-$\tfrac12$ observables $U_i=\mathbf n_i\cdot\boldsymbol\sigma$,
\begin{equation}\label{eq:spinweight}
c_{ij}=|\mathbf n_i\times\mathbf n_j|=|\sin\theta_{ij}|,
\qquad
D_{\mathrm{spin}}(\sigma,\tau)=\max_{\Inv(\sigma,\tau)}|\sin\theta_{ij}|.
\end{equation}
